\documentclass[25pt, a0paper, landscape, margin=25.4mm]{tikzposter}

\usepackage{amsmath}
\usepackage{amssymb}
\usepackage[backend=bibtex]{biblatex}
\usepackage{filecontents}

\bibliography{bibliography}
\bibstyle{apalike}

\newcommand{\T}[1]{\texttt{#1}}

\makeatletter
\def\title#1{\gdef\@title{\scalebox{\TP@titletextscale}{%
\begin{minipage}[t]{\linewidth}
\centering
#1
\par
\vspace{0.5em}
\end{minipage}%
}}}
\makeatother

\makeatletter
\renewenvironment{tikzfigure}[1][]{
  \def \rememberparameter{#1}
  \vspace{10pt}
  \refstepcounter{figurecounter}
  \begin{center}
  }{
    \ifx\rememberparameter\@empty
    \else %nothing
    \\[10pt]
    {\Large Fig.~\thefigurecounter: \rememberparameter}
    \fi
  \end{center}
}
\makeatother

\title{Extraction of Cost Recurrences from Sequential and Parallel Functional Programs}
\institute{Wesleyan University} % See Section 4.1
\author{Justin Raymond}
%\titlegraphic{\includegraphics[height=50pt]{wesleyan_shield.png}}

\usetheme{Simple}  % See Section 5

\begin{document}
\maketitle  % See Section 4.1

\begin{columns}  % See Section 4.4

\column{0.33}
%
% ABSTRACT
%
\block{Abstract}{
  \Large
Complexity analysis aims to predict the resources, most often time and space,
which a program requires. Traditional complexity analysis is not compositional,
often does not address higher order functions, and there is no formal relation
between programs and the equations that describe their complexity. We build on
previous work by \cite{Danner2013} and \cite{Danner2015} which formalizes the
extraction of cost recurrences from higher order functional programs. We use
the formalization to analyze the time complexity of higher order functional
programs. We also demonstrate the flexibility of the method by extending it
to parallel cost semantics.
}
\block{Complexity Analysis}{
\LARGE
The higher order function \T{fold} reduces a list to a single element using a
combining function.
%
\begin{align*}
  \T{fold} &= \lambda f\ z\ xs.\T{rec}(xs,\T{Nil} \mapsto z, \\
           &\qquad \T{Cons}\mapsto (x,xs',r).f\ x\ \T{force}(r))
\end{align*}
%
\normalfont
To analyze the complexity of \T{fold}, we write down a recurrence $T(n)$ which
reflects the number of steps to execute the \T{fold} applied to a function $f$,
a seed value $z$, and a list $xs$.
%
\[ T(n) = c_f + T(n-1) \]
%
The closed form solution is $T(n) = c_f\ n$. The method we used to write the
recurrence for the program was informal. The analysis is not composable: we
cannot use our results to analyze \T{map f $\circ$ fold g z}. The analysis also
assumes the cost of applying $f$ to each item in the list is constant.
}
%
% Higher Order Complexity Analysis
%
\column{0.33}
\block{Higher Order Complexity Analysis}{
  \LARGE
  Instead of considering the complexity of a program as a cost, we consider it
  as a pair of a cost and a potential.  The cost is a bound on the steps
  required to run the program. The potential is a measure of the cost of future
  use of the program.
  %
  \begin{tikzfigure}[The type of the translation function.]
    $\|\tau\| = \textbf{C} \times \langle\!\langle \tau \rangle\!\rangle$
  \end{tikzfigure}
  %
  To formally extract a recurrence from a program, we translate it from its
  original "source" language into a "complexity" language. We then interpret
  the complexity language program in a denotational semantics to obtain the
  recurrence.
  %
  \begin{tikzfigure}[The translation of a tuple.]
    $\|\langle e_0,e_1\rangle\| = \langle \|e_0\|_c + \|e_1\|_c, \langle \|e_0\|_p \|e_1\|_p\rangle\rangle$
  \end{tikzfigure}
  %
  This is a formal process for extracting a recurrence for the cost of executing
  a program, and the recurrence is a bound on the cost of executing the program.

  Since the analysis of a program results in a measure of future use of the
  program, we can compose analysis of programs to analyze the composition of
  the programs as well analyze higher-order programs.
}
\block{References}{
  \printbibliography[heading=none]
}
%
% Parallel Complexity Analysis
%
\column{0.33}
\block{Parallel Complexity Analysis}{
  \LARGE
  To analyze the complexity of a parallel program we use a specification that is
  agnostic to the details of how to schedule subcomputations, that is the
  analysis does not depend on the number of processors available. Instead of
  representing costs as natural numbers, we represent a cost as a graph which
  describes the dependencies between subcomputations of the program.
  %
  \begin{tikzfigure}[Cost graph definition]
    $\textbf{C} = 0\ |\ 1\ |\ \textbf{C} \oplus \textbf{C}\ |\ \textbf{C} \otimes \textbf{C}$
  \end{tikzfigure}

  Two measures can be extracted from the cost
  graph: work, the total amount of steps to execute a program, and span, the length
  of the critical path of the program execution. The predicted running time of
  a program with work $w$ and span $s$ on $p$ processors is given by Brent's
  Theorem:
  \[
    O(max(\frac{w}{p}, s))
  \]
  %
  The translation rules must be altered to produce cost graphs instead of
  natural numbers.
  %
  \begin{tikzfigure}[Parallel Translation of a tuple]
      $\|\langle e_0,e_1\rangle\| = \langle \|e_0\|_c \otimes \|e_1\|_c, \langle \|e_0\|_p,\|e_1\|_p\rangle\rangle$
  \end{tikzfigure}
  %
  The steps to execute the source language program is bounded by the
  recurrence in the translated complexity language program.
  %
  \begin{tikzfigure}[Bounding Theorem]
    If $\gamma \vdash e : \tau$, then $e \sqsubseteq \|e\|$
  \end{tikzfigure}
  %
  $\sqsubseteq$ is a binary  logical relation stating the cost of executing
  the program $e$ is bounded by the cost of $\|e\|$.
}

\end{columns}


\end{document}
